\documentclass[a4paper,10pt]{amsart}

\pdfinfoomitdate=1
\pdftrailerid{}
\pdfsuppressptexinfo=15

\usepackage[T1]{fontenc}
\usepackage{lmodern}
\usepackage[margin=27mm]{geometry}
\usepackage{amsmath,amssymb,mathtools,booktabs,array,microtype,xcolor}
\usepackage[shortlabels]{enumitem}
\usepackage[numbers,sort&compress]{natbib}
\usepackage[colorlinks=true,linkcolor=blue!55!black,citecolor=blue!55!black,urlcolor=blue!55!black]{hyperref}
\hypersetup{pdftitle={},pdfauthor={},pdfsubject={},pdfkeywords={}}

\newtheorem{theorem}{Theorem}
\newtheorem{proposition}[theorem]{Proposition}
\newtheorem{lemma}[theorem]{Lemma}
\newtheorem{corollary}[theorem]{Corollary}
\theoremstyle{remark}
\newtheorem{remark}[theorem]{Remark}

\newcommand{\Prob}{\mathbb P}
\newcommand{\E}{\mathbb E}
\newcommand{\Br}{B_r}

\title[Triad dynamics on triangular books]{Joint Absorption and Spine-Flip Laws for Local Triad Dynamics on Triangular Books}
\author{Anonymous}
\date{}

\begin{document}

\begin{abstract}
For the established $p=1/3$ local triad dynamics, we specialize to the
triangular book $B(3,r)=K_{1,1,r}$ and count active imbalanced-triad update
epochs.  The page-imbalance count is a reflected one-dimensional quotient.
We solve its transform jointly marked by absorption time and common-spine
flips in an explicit Chebyshev-rational form.  The mean absorption time is the
quadratic $k(r+2-k)/2$, with sharp extrema over the starting count.  The mean
together with the probability of odd spine parity exactly identifies the
book size and initial imbalance count, subject to an explicit integer
feasibility criterion, while either statistic alone has concrete collisions.
A private-edge block supplies a uniform absorption certificate.  All formulas
include the $r=1$, $r=2$, $z=0$, coincident-arrow, and clock-convention
boundaries.  An exact-arithmetic audit performs $199{,}581$ rational and
integer checks.
\end{abstract}

\maketitle

\section{Owned process, special carrier, and theorem}\label{sec:setup}

Let
\[
 \Br=B(3,r)=K_{1,1,r},\qquad r\geq1,
\]
be the graph obtained by joining $r$ triangles along one common \emph{edge},
called the spine.  Thus each page has two private edges.  Give every physical
edge a sign in $\{\pm1\}$, and call a page imbalanced when the product of its
three edge signs is negative.  At an update epoch, choose a currently
imbalanced page uniformly, choose one of its three physical edges uniformly,
and flip that sign.  Stop when all pages are balanced.

Write $x_i\in\{0,1\}$ for the imbalance bit of page $i$ and
$K=\sum_i x_i$.  Let $T$ be the number of these \emph{active update epochs}
until absorption, and let $J$ be the number of spine flips before absorption.
Expectations conditional on any sign state with $K=k$ are denoted by $\E_k$;
strong lumpability below makes this notation unambiguous.

This stochastic rule is the $p=1/3$ specialization of local triad dynamics
of Antal--Krapivsky--Redner~\cite{AntalKrapivskyRedner2005,AntalKrapivskyRedner2006}.
The 2005 all-triad clock samples balanced pages too and hence inserts no-op
holds; the process above is its chain embedded at active epochs.  The 2006
formulation directly selects an imbalanced triad.  Istrate owns the same
probabilistic kernel on general graphs and its triadic-dual/XOR
representation~\cite{Istrate2009}; later hypergraph particle-system and drift
methods are also background~\cite{IstrateBonchisMarin2019}.  Signed-book work
owns the carrier and its static switching-class count indexed by negative
pages~\cite{SehrawatBhattacharjya2022}.  Table~\ref{tab:subtraction} makes the
resulting claim boundary explicit.

\begin{table}[ht]
\centering
\small
\caption{Primary-source subtraction.  The middle column receives zero
contribution credit; the right column is the special-carrier residual.}
\label{tab:subtraction}
\begin{tabular}{@{}>{\raggedright\arraybackslash}p{.21\textwidth}
                    >{\raggedright\arraybackslash}p{.35\textwidth}
                    >{\raggedright\arraybackslash}p{.34\textwidth}@{}}
\toprule
Source & Owned input & Residual used here \\
\midrule
Antal--Krapivsky--Redner~\cite{AntalKrapivskyRedner2005,AntalKrapivskyRedner2006}
& local triad dynamics and the $p=1/3$ equiprobable-edge update
& exact joint law after specialization to $B(3,r)$ and the active clock \\
Istrate and Istrate--Bonchis--Marin~\cite{Istrate2009,IstrateBonchisMarin2019}
& same general-graph kernel, triadic dual/XOR encoding, hypergraph
particle systems, drift, and convergence-time program
& reflected count recurrence, closed transform, and coarse inverse on the
one-spine carrier only \\
Sehrawat--Bhattacharjya~\cite{SehrawatBhattacharjya2022}
& signed-book carrier and static switching classes indexed by page count
& no carrier, switching, coloring, or static-class claim \\
Generic finite-chain tools
& Bellman/resolvent rationality, Chebyshev facts, tail sums, and quadratic
concavity
& only the boundary-complete formula conjunction stated below \\
\bottomrule
\end{tabular}
\end{table}

For $0\leq k\leq r$, define the joint transform
\begin{equation}\label{eq:F-def}
 F_k(z,u)=\E_k[z^T u^J],\qquad F_0(z,u)=1.
\end{equation}
Let $U_j$ denote the Chebyshev polynomial of the second kind, with
$U_{-1}=0$, $U_0=1$, and
$U_{j+1}(x)=2xU_j(x)-U_{j-1}(x)$.

\begin{theorem}[joint law, clock, inverse, and absorption]\label{thm:main}
For the process on $\Br$ started with $K=k$, $1\leq k\leq r$, the following
hold.
\begin{enumerate}[(i)]
\item The full sign chain is strongly lumpable by $K$, with count transitions
\begin{equation}\label{eq:quotient}
 k\longrightarrow k-1\quad\text{with probability }\frac23,
 \qquad
 k\longrightarrow r-k\quad\text{with probability }\frac13.
\end{equation}
If the targets coincide, the two masses add.

\item For $r\geq2$ and $z\neq0$, put
\begin{equation}\label{eq:xi}
 \xi=\frac{9+z^2(4-u^2)}{12z}.
\end{equation}
Then, as identities of reduced rational functions,
\begin{align}
 F_k(z,u)&=U_{k-1}(\xi)F_1(z,u)-U_{k-2}(\xi),
       \label{eq:Fk}\\
 F_1(z,u)&=
 \frac{3U_{r-2}(\xi)-2zU_{r-3}(\xi)+zu}
      {3U_{r-1}(\xi)-2zU_{r-2}(\xi)}.             \label{eq:F1}
\end{align}
For probability-transform evaluation, these identities hold on
$|z|<1$, $|u|\leq1$, with removable points taken from the Bellman system.
The boundary values are
\begin{equation}\label{eq:small-boundaries}
 r=1:\quad F_1=\frac{z(2+u)}3;\qquad
 r=2:\quad F_1=\frac{2z}{3-zu};\qquad
 z=0:\quad F_k=\mathbf 1_{\{k=0\}}.
\end{equation}

\item The mean is
\begin{equation}\label{eq:mean}
 m_k:=\E_kT=\frac{k(r+2-k)}2.
\end{equation}
For $r>1$, its unique minimum over nonabsorbing counts is at $k=1$,
with value $(r+1)/2$.  Its maxima are
\begin{equation}\label{eq:extrema}
 \begin{cases}
 k=(r+2)/2,\quad m_k=(r+2)^2/8, & r\text{ even},\\
 k=(r+1)/2,(r+3)/2,\quad m_k=((r+2)^2-1)/8,
   & r\text{ odd}.
 \end{cases}
\end{equation}
For $r=1$, the sole nonabsorbing state is both extremizers and has mean one.

\item If $q_k=\Prob_k(J\text{ is odd})$, then
\begin{equation}\label{eq:parity}
 \E_k[(-1)^J]=\frac{r+2-2k}{r+2},
 \qquad q_k=\frac{k}{r+2}.
\end{equation}
Given an exact real candidate pair $(m,q)$ for a nonabsorbing start, declare
it infeasible unless $m>0$ and $0<q<1$.  On that domain set
\begin{equation}\label{eq:R}
 R=\sqrt{\frac{2m}{q(1-q)}}.
\end{equation}
The pair is feasible if and only if $R$ is an integer with $R\geq3$ and
$k=qR$ is an integer satisfying $1\leq k\leq R-2$.  In that case
\begin{equation}\label{eq:inverse}
 r=R-2,\qquad k=qR
\end{equation}
are uniquely recovered.  Both observations are necessary in general.

\item For every integer $n\geq0$,
\begin{equation}\label{eq:tail}
 \Prob_k(T>nr)\leq\left[1-\left(\frac23\right)^r\right]^n.
\end{equation}
In particular, absorption is almost sure and $T$ has an exponential tail.
\end{enumerate}
\end{theorem}

The theorem recovers only the carrier size and initial imbalance \emph{count}.
It neither reconstructs a full edge-sign configuration nor supplies stability
under noisy observations.

\section{Strong lumping and an absorption certificate}\label{sec:lumping}

\begin{proof}[Proof of Theorem~\ref{thm:main}(i) and (v)]
Flipping either private edge of the selected active page clears that page and
changes no other imbalance bit.  Both private choices therefore send every
state of count $k$ to the count $k-1$.  Flipping the common spine toggles all
$r$ imbalance bits and sends count $k$ to $r-k$.  The probabilities are the
same for every full sign state within the count fibre, which proves strong
lumpability and \eqref{eq:quotient}.  If $k-1=r-k$, the private and spine
choices are different physical updates with the same quotient target, so
their probabilities combine rather than producing two arrows.

Pre-generate the private/spine type of the next $r$ edge choices.  At each
active epoch the type is private with probability $2/3$, independently of the
state.  If all $r$ types are private, each performed update lowers $K$ by one,
so absorption occurs within the block, possibly before all pre-generated
choices are used.  This event has probability $(2/3)^r$.  Conditional on
survival to any block boundary, the same bound applies.  Iteration with the
Markov property proves \eqref{eq:tail}, and hence almost-sure absorption.
\end{proof}

On the friendship or Dutch-windmill graph, the triangles share only a vertex,
not an edge.  Every physical edge is then private, so $K\mapsto K-1$ and
$T=K$ deterministically.  That carrier does not satisfy
\eqref{eq:quotient}.

\section{The spine-marked Chebyshev transform}\label{sec:transform}

First-step conditioning in the quotient gives, for $1\leq k\leq r$,
\begin{equation}\label{eq:bellman}
 F_k=z\left(\frac23F_{k-1}+\frac{u}{3}F_{r-k}\right).
\end{equation}
The factor $u$ marks precisely a common-spine flip; it does not mark all sign
flips or the terminal number of negative edges.

\begin{proof}[Proof of Theorem~\ref{thm:main}(ii)]
Work first in the rational-function field with $z\neq0$.  For
$1\leq k<r$, equations \eqref{eq:bellman} at $k$ and $r-k$ give
\begin{equation}\label{eq:elim-pairs}
 uF_{r-k}=\frac{3}{z}F_k-2F_{k-1},\qquad
 F_{r-k-1}=\frac{3}{2z}F_{r-k}-\frac{u}{2}F_k.
\end{equation}
Insert the second identity into \eqref{eq:bellman} at $k+1$, then use the
first.  No division by $u$ is needed, and simplification gives
\begin{equation}\label{eq:cheb-rec}
 F_{k+1}=\frac{9+z^2(4-u^2)}{6z}F_k-F_{k-1}
        =2\xi F_k-F_{k-1}.
\end{equation}
Since $F_0=1$, the Chebyshev solution of \eqref{eq:cheb-rec} is
\eqref{eq:Fk}.  At the terminal count, \eqref{eq:bellman} becomes
\begin{equation}\label{eq:terminal}
 F_r=\frac{2z}{3}F_{r-1}+\frac{zu}{3}.
\end{equation}
Substituting \eqref{eq:Fk} into \eqref{eq:terminal} and collecting the
$F_1$ terms gives \eqref{eq:F1}.

For $r=1$, every edge choice absorbs in one epoch; two choices are private
and one is the marked spine, which gives the first formula in
\eqref{eq:small-boundaries}.  When $r=2$, a spine flip from $k=1$ is a
self-loop, so \eqref{eq:bellman} directly gives $F_1=2z/(3-zu)$.  Equivalently,
the unreduced Chebyshev ratio is
\[
 \frac{3+zu}{6\xi-2z},\qquad
 6\xi-2z=\frac{(3-zu)(3+zu)}{2z};
\]
the removable factor $3+zu$ must be cancelled before pointwise evaluation.
Finally, $\xi$ is undefined at $z=0$, whereas \eqref{eq:bellman} gives
$F_0=1$ and $F_k=0$ for $k>0$.  Thus \eqref{eq:Fk}--\eqref{eq:F1} are read as
reduced rational identities with these Bellman continuations.
\end{proof}

Uniform choices matter.  Nonuniform page weights generally destroy count
lumpability, while nonuniform physical-edge weights change the coefficients
in both \eqref{eq:quotient} and \eqref{eq:bellman}.

\section{Quadratic clock and exact coarse inverse}\label{sec:inverse}

\begin{proof}[Proof of Theorem~\ref{thm:main}(iii)]
The exponential tail justifies first-step conditioning of the mean:
\begin{equation}\label{eq:mean-bellman}
 m_0=0,\qquad
 m_k=1+\frac23m_{k-1}+\frac13m_{r-k}.
\end{equation}
Use \eqref{eq:mean-bellman} at $k$, $r-k$, and $k+1$.  Eliminating the two
reflected terms yields
\[
 m_{k+1}-2m_k+m_{k-1}=-1\qquad(1\leq k<r).
\]
The terminal equation is $m_r=1+(2/3)m_{r-1}$.  Together with $m_0=0$, these
conditions have the unique quadratic solution \eqref{eq:mean}.

For $r>1$, strict concavity makes the minimum occur at an endpoint; since
$m_1=(r+1)/2<m_r=r$, it is uniquely at $k=1$.  The vertex is
$(r+2)/2$, so the nearest admissible integer or integers give exactly
\eqref{eq:extrema}.  The one-state case $r=1$ is immediate.
\end{proof}

\begin{proof}[Proof of Theorem~\ref{thm:main}(iv)]
Put $h_k=\E_k[(-1)^J]$.  A private flip preserves the parity sign and a spine
flip reverses it, whence
\begin{equation}\label{eq:parity-bellman}
 h_0=1,\qquad h_k=\frac23h_{k-1}-\frac13h_{r-k}.
\end{equation}
The affine values $(r+2-2k)/(r+2)$ satisfy every equation.  Iterating the
bounded first-step identity to $T\wedge n$ and using almost-sure absorption
and bounded convergence identifies them with $\E_k[(-1)^J]$.  The second
identity in \eqref{eq:parity} follows from
$\E[(-1)^J]=1-2\Prob(J\text{ odd})$.

Let $R_0=r+2$.  Equations \eqref{eq:mean} and \eqref{eq:parity} give
\[
 m=\frac{k(R_0-k)}2,\qquad
 q(1-q)=\frac{k(R_0-k)}{R_0^2}.
\]
In particular, every genuine pair has $m>0$ and $0<q<1$.  Therefore
\eqref{eq:R} equals the positive integer $R_0$, and $qR_0=k$ is an
integer in $[1,R_0-2]$.  Conversely, if the stated feasibility conditions
hold, the pair $r=R-2$, $k=qR$ is admissible and the two displayed identities
reproduce $(m,q)$.  This proves both feasibility and uniqueness.

The central value $q=1/2$ is regular: it occurs when $r$ is even and
$k=(r+2)/2$.  Neither scalar alone suffices.  Indeed,
\[
 (r,k)=(1,1),(4,2)\quad\text{both give }q=\frac13,
 \qquad
 (r,k)=(2,2),(3,1)\quad\text{both give }m=2.
\]
\end{proof}

\section{Exact pressure, controls, and limitations}\label{sec:audit}

A standard-library verifier enumerates every nonzero imbalance vector through
$r=9$, solves independent rational Bellman systems, and compares them with the
closed formulas.  It also compares the inverse criterion with the literal
two-statistic image on $7{,}335$ bounded exact candidate pairs, including
$7{,}266$ rejections, checks both printed scalar collisions, enumerates exact
private-block masses, and propagates $546$ finite tail probabilities.
Table~\ref{tab:audit} records the frozen lanes.

\begin{table}[ht]
\centering
\small
\caption{Deterministic exact-arithmetic falsification.  Enumeration supports
debugging; the preceding symbolic arguments prove the all-parameter theorem.}
\label{tab:audit}
\begin{tabular}{@{}lr@{}}
\toprule
Lane & Exact assertions \\
\midrule
Chebyshev elimination and transform vectors & 4,416 \\
Inverse iff grid and scalar collisions & 7,655 \\
Inverse states and absorption certificate & 180,600 \\
Literal strong lumpability & 2,026 \\
Mean, parity, and extrema & 3,958 \\
Private-block probability and tail & 648 \\
$r=1$, $r=2$, and $z=0$ boundaries & 278 \\
\midrule
Total & 199,581 \\
\bottomrule
\end{tabular}
\end{table}

The exact control uses only integers and rational numbers.  It is neither a
proof of the theorem nor an ownership, novelty, priority, or release
certificate.  The primary-source screen is bounded, and any direct owner of
the special-book law would require claim subtraction anew.

Internally adjacent objects are separated by the literal update.  A
triangular-book edge-deletion process removes edges; here every edge remains
and the spine flip reflects $k$ to $r-k$.  Vertex-push orientation chains are
stationary group walks, whereas this process selects an active imbalanced
page and absorbs.  Deterministic prefix-XOR feedback and unequal-spider first
passage use different state spaces and proof engines; the bivariate
spine-marked reflection law, rather than XOR notation or a mean alone, is the
organizing object here.

\section*{Limitations}
The results require uniform selection of active pages and physical edges,
use the active update-epoch clock, and assume the known carrier
$B(3,r)=K_{1,1,r}$.  They do not cover the all-triad physical clock,
nonuniform weights, arbitrary graphs, friendship/windmill carriers, noisy
inverse data, or recovery of a full sign state.  The inverse uses exact
$(m,q)$ and identifies only $(r,k)$.

\section*{Data Availability}
No external data were used.  The paper-local verifier and its frozen
transcript contain the complete deterministic exact control.

\section*{Ethics Statement}
This mathematical study involved no human participants, animals, personal
data, or field intervention.

\section*{Author Contributions}
The anonymous author performed the mathematical derivations, exact checks,
source-boundary audit, and manuscript preparation.

\section*{Conflict of Interest}
The author declares no conflict of interest.

\section*{Funding}
No external funding is declared.

\section*{External Status}
This artifact remains \texttt{HOLD\_EXTERNAL}; it is not cleared for posting,
submission, circulation, or author contact.

\bibliographystyle{plainnat}
\bibliography{references}

\end{document}
